\documentclass[a4paper,11pt]{ltjsarticle}
\usepackage{luatexja}
\usepackage{luatexja-fontspec}
\usepackage{amsmath,amssymb,amsthm}
\usepackage{mathtools}
\usepackage{booktabs}
\usepackage{longtable}
\usepackage{array}
\usepackage{hyperref}
\usepackage{graphicx}
\usepackage{float}
\usepackage{geometry}
\geometry{margin=25mm}

\hypersetup{
  colorlinks=true,
  linkcolor=blue,
  urlcolor=blue,
  citecolor=blue
}

\theoremstyle{definition}
\newtheorem{theorem}{定理}[section]
\newtheorem{proposition}[theorem]{命題}
\newtheorem{lemma}[theorem]{補題}
\newtheorem{corollary}[theorem]{系}
\newtheorem{definition}[theorem]{定義}
\newtheorem{remark}[theorem]{注意}
\newtheorem{observation}[theorem]{観察}

\setlength{\parindent}{1em}
\setlength{\parskip}{0.5em}

\begin{document}

\begin{center}
{\LARGE\bfseries 複数の主観空間における観測の相対性：\\中心投影の対称性の幾何学的帰結}
\end{center}

\bigskip

\noindent\textbf{著者：} Noriaki Kihara（木原 範昭）\\
\textbf{所属：} WF System Co., Ltd.（大阪大学 基礎工学部 卒業）\\
\textbf{作成日：} 2026年4月\\
\textbf{種別：} 研究ノート（幾何学的考察）\\
\textbf{前論文：}\\
{[1]} 中心投影による4次元空間の幾何学的定式化. DOI: \href{https://doi.org/10.5281/zenodo.19427780}{10.5281/zenodo.19427780}.\\
{[2]} 中心投影の幾何学的対称性：多軸モデルの数学的基盤. DOI: \href{https://doi.org/10.5281/zenodo.19434932}{10.5281/zenodo.19434932}.

\bigskip\hrule\bigskip

\section*{概要}

前2稿 {[1, 2]} において，中心投影（central projection）による \(n\) 次元接超平面から超球面 \(S^n(R)\) への写像を構成し，その4つの幾何学的対称性を証明した．本稿では，論文 {[2]} の対称性 II（軸の対等性）と対称性 IV（主観座標系の変換可能性）の自然な帰結として，\textbf{異なる軸を投影中心とする複数の主観空間が同じ背景空間から構成可能であり，それぞれの主観空間における観測量は内部観測者の立場の選択に依存する}ことを示す．

具体的には：

\begin{enumerate}
\item 背景空間 \(\mathbb{R}^{n+1}\) から，投影中心軸の選択により異なる主観空間 \(H_A^+, H_B^+, \ldots\) が構成される
\item 各主観空間における内部観測者は，計量から構成される量のみを直接観測する
\item 異なる主観空間の内部観測者は，同じ背景空間を共有しながら，異なる観測量を取得する
\item 論文 {[2]} 対称性 IV により，異なる主観空間の間には背景空間を経由した正則な合成写像が存在する
\end{enumerate}

これは中心投影の対称性 II と IV から直接導かれる純粋に幾何学的な帰結であり，物理的解釈を一切含まない．本論文は新たな仮定を導入せず，論文 {[1, 2]} の結果のみを用いて論証する．

\bigskip\hrule\bigskip

\section{はじめに}

論文 {[1]} は，\(n\) 次元接超平面 \(\Pi_R\) から超球面 \(S^n(R) \subset \mathbb{R}^{n+1}\) への中心投影を定義し，誘導計量と Einstein テンソルを導出した．論文 {[2]} は，この中心投影が持つ4つの幾何学的対称性を証明した：

\begin{itemize}
\item \textbf{対称性 I}（離散安定性）：座標が零・正の整数・負の整数のいずれの値をとっても写像は正則
\item \textbf{対称性 II}（軸の対等性）：背景空間 \(\mathbb{R}^{n+1}\) のどの軸を投影中心に選んでも写像の構造は同一
\item \textbf{対称性 III}（測地線偏差）：曲率半径 \(R\) が測地線偏差 \(|\xi|^2\) を生む
\item \textbf{対称性 IV}（主観座標系の変換可能性）：異なる軸を中心とする2つの主観座標系は相互変換可能
\end{itemize}

本稿では，これらの対称性のうち特に対称性 II と対称性 IV から自然に導かれる構造を述べる．具体的には，背景空間 \(\mathbb{R}^{n+1}\) から異なる軸を投影中心として構成される複数の主観空間が同時に存在しうること，そしてそれぞれの主観空間における内部観測者は異なる観測量を取得することを示す．

本論文の主張はすべて純粋に幾何学的な命題であり，物理的解釈を一切含まない．本論文の主張範囲は \S7 で明示的に限定される．

\section{準備：複数の主観空間}

\subsection{中心投影と主観空間の記法}

論文 {[2]} \S2.1 の定義に従い，背景空間 \(\mathbb{R}^{n+1}\) における基底ベクトル \(e_A\)（\(A \in \{1, \ldots, n+1\}\)）を投影中心軸とする中心投影 \(\Phi_A\) を考える：
\begin{equation}
\Phi_A: \Pi_R^{(A)} \to S^n(R), \quad Y^\mu = \frac{Rx^\mu}{l_A} \quad (\mu \neq A), \qquad Y^A = \frac{R^2}{l_A} \tag{2.1}
\end{equation}
ここで \(l_A = \sqrt{R^2 + \sum_{\mu \neq A}(x^\mu)^2}\) である．\(\Phi_A\) の像は開半球面
\begin{equation}
H_A^+ = \{Y \in S^n(R) \mid Y^A > 0\} \tag{2.2}
\end{equation}
であり，これを\textbf{軸 \(A\) を投影中心とする主観空間}と呼ぶ．

\subsection{複数の主観空間の同時存在}

論文 {[2]} 対称性 II（定理 4.1，軸の対等性）により，背景空間 \(\mathbb{R}^{n+1}\) のどの軸 \(A\) を投影中心に選んでも，中心投影 \(\Phi_A\) は同一の代数的構造を持ち，引き戻し計量・曲率テンソル・Einstein テンソルはいずれも軸の選択 \(A\) に依存しない．

したがって，背景空間 \(\mathbb{R}^{n+1}\) から \(n+1\) 通りの主観空間 \(\{H_A^+\}_{A=1}^{n+1}\) が同時に構成可能である．これら \(n+1\) 個の主観空間は，同一の背景空間を共有しながら，それぞれ異なる軸を投影中心とする．

\begin{definition}[主観空間の族]
背景空間 \(\mathbb{R}^{n+1}\) から構成される主観空間の族を
\begin{equation}
\mathcal{H} = \{H_A^+ \mid A \in \{1, 2, \ldots, n+1\}\} \tag{2.3}
\end{equation}
と表す．
\end{definition}

\begin{proposition}[主観空間の構造的同一性]
任意の \(A, B \in \{1, \ldots, n+1\}\) に対し，主観空間 \(H_A^+\) と \(H_B^+\) は構造的に同一の幾何学的性質を持つ：

\textup{(i)} 引き戻し計量はそれぞれの主観座標で同一の代数的形を取る

\textup{(ii)} リーマン曲率テンソル，リッチテンソル，Einstein テンソル，スカラー曲率はすべて同一の値を取る

\textup{(iii)} 断面曲率は \(K = 1/R^2\) で共通
\end{proposition}

\begin{proof}
論文 {[2]} 定理 4.1 (ii) より直接従う．軸の対等性により，引き戻し計量は
\begin{equation}
g_{\mu\nu}^{(A)} = \frac{R^2}{l_A^2}\left(\delta_{\mu\nu} - \frac{x_\mu x_\nu}{l_A^2}\right) \quad (\mu, \nu \neq A) \tag{2.4}
\end{equation}
の形を取り，これは \(A\) の選択に依存しない代数的構造である．曲率テンソル類はすべて計量から決まり，定曲率空間（{[1]} 定理 4.1）の性質として軸の選択に依存しない．
\end{proof}

\subsection{内部観測者と主観空間}

論文 {[2]} \S2 用語表に従い，\textbf{内部観測者}を主観空間 \(H_A^+\) の内部で計量 \(g_{\mu\nu}^{(A)}\) のみを用いて測定を行う仮想的な観測者と定義する．以下，主観空間 \(H_A^+\) の内部観測者を\textbf{観測者 \(A\)} と呼ぶ．

観測者 \(A\) は主観空間 \(H_A^+\) の主観座標 \((x^\mu)_{\mu \neq A}\) を用いて観測を記述する．論文 {[2]} 定理 5.2 (i) により，曲率半径 \(R\) は観測者 \(A\) には \([L]\) として直接測定不可能であり，\(R^2\) を通じてのみ計量に寄与する．観測者 \(A\) が取得できる情報は計量から構成されるスカラー不変量に限定される．

\section{主観空間内の観測}

\subsection{観測者 \(A\) の取得情報}

主観空間 \(H_A^+\) の内部観測者 \(A\) が論文 {[1, 2]} の枠組みで取得できる情報を整理する．

\begin{proposition}[観測者 \(A\) の取得情報]
観測者 \(A\) は次の情報を取得する：

\textup{(i)} 主観座標 \((x^\mu)_{\mu \neq A}\) における自分の位置と運動

\textup{(ii)} 計量 \(g_{\mu\nu}^{(A)}\) から構成されるスカラー不変量（測地線偏差 \(|\xi|^2\)，スカラー曲率 \(n(n-1)/R^2\) など）

\medskip
観測者 \(A\) が取得\textbf{しない}情報：

\textup{(iii)} 投影中心軸 \(A\) の値（論文 {[2]} 定理 5.2 (i) により \([L]\) として測定不可能）

\textup{(iv)} 背景空間 \(\mathbb{R}^{n+1}\) における絶対位置（論文 {[2]} 命題 6.5：位置情報の構造的脱落）

\textup{(v)} 主観空間 \(H_A^+\) の外側にある他の主観空間 \(H_B^+\)（\(B \neq A\)）の存在や構造（観測者 \(A\) は自分の主観空間 \(H_A^+\) の内部にのみアクセスする）
\end{proposition}

\begin{proof}
(i)--(ii) は論文 {[1, 2]} の結果から直接従う．(iii) は論文 {[2]} 定理 5.2 (i) による．(iv) は論文 {[2]} 命題 6.5 による．(v) は内部観測者の定義（論文 {[2]} \S2 用語表）から従う：内部観測者は自分の住む主観空間の計量のみを用いて測定を行う仮想的な観測者であり，他の主観空間にはアクセスしない．
\end{proof}

\subsection{観測者 \(A\) の主観空間における軸の役割}

観測者 \(A\) は主観空間 \(H_A^+\) の中で，自分の主観座標 \((x^\mu)_{\mu \neq A}\) の各軸について観測を行う．これら \(n\) 個の軸は，観測者 \(A\) にとって「自分が住む空間の軸」として現れる．

ここで重要な点を述べる．観測者 \(A\) は，自分の主観空間が「軸 \(A\) を投影中心として構成された」ことを知らない．観測者 \(A\) にとって，自分の主観座標 \((x^\mu)_{\mu \neq A}\) は単に「自分の空間の座標」であり，どの軸が背景空間における投影中心軸であったかという情報は，命題 3.1 (iii) により取得できない．

観測者 \(A\) にとっての主観空間 \(H_A^+\) は，自己完結した世界として現れる．背景空間 \(\mathbb{R}^{n+1}\) の存在も，他の主観空間 \(H_B^+\) の存在も，観測者 \(A\) には直接アクセスできない情報である．

\section{異なる主観空間の対比}

\subsection{観測者 \(A\) と観測者 \(B\) の対比}

異なる二つの主観空間 \(H_A^+\) と \(H_B^+\)（\(A \neq B\)）における内部観測者，観測者 \(A\) と観測者 \(B\) を対比する．両者は同じ背景空間 \(\mathbb{R}^{n+1}\) から構成された主観空間に住むが，それぞれの観測する量は異なる役割を持つ．

\begin{proposition}[観測者 \(A\) と観測者 \(B\) の対比]
観測者 \(A\) と観測者 \(B\) について，次の対比が成立する：

\textup{(i)} 観測者 \(A\) は主観座標 \((x^\mu)_{\mu \neq A}\) を観測し，観測者 \(B\) は主観座標 \((x'^\nu)_{\nu \neq B}\) を観測する．両者の主観座標は，論文 {[2]} 定理 6.1 により背景空間 \(\mathbb{R}^{n+1}\) を経由した合成写像 \(T_{A \to B} = \Phi_B^{-1} \circ \Phi_A\) で関係づけられる．

\textup{(ii)} 観測者 \(A\) にとって軸 \(A\) は「主観空間の外側」（投影中心軸）であり直接アクセスできない．一方，軸 \(A\) は観測者 \(B\) の主観座標の一つである（\(A \neq B\) であるため，軸 \(A\) は観測者 \(B\) の主観空間 \(H_B^+\) 内の座標 \(x'^A\) として現れる）．

\textup{(iii)} 同様に，軸 \(B\) は観測者 \(B\) にとって「主観空間の外側」であり直接アクセスできないが，観測者 \(A\) の主観座標の一つ \(x^B\) として現れる．
\end{proposition}

\begin{proof}
(i) 論文 {[2]} 定理 6.1 により，\(\Phi_A\) と \(\Phi_B\) の合成 \(T_{A \to B} = \Phi_B^{-1} \circ \Phi_A\) は \(C^\infty\) 級の微分同相写像である．観測者 \(A\) の主観座標と観測者 \(B\) の主観座標は，この合成写像を通じて関係づけられる．

(ii)--(iii) 論文 {[2]} \S2.1 より，主観座標 \(\Phi_A^{-1}\) により得られる座標は \((x^\mu)_{\mu \neq A}\) であり，軸 \(A\) そのものは含まれない．軸 \(A\) は背景空間 \(\mathbb{R}^{n+1}\) の基底ベクトルであり，観測者 \(A\) の主観空間 \(H_A^+\) には現れない．一方，観測者 \(B\) の主観空間 \(H_B^+\) では \(A \neq B\) である限り，背景空間の軸 \(A\) は主観座標 \(x'^A\) として現れる．論文 {[2]} 命題 6.1 の式 (6.5) により，観測者 \(B\) の主観座標における \(x'^A\) は観測者 \(A\) の主観座標における \(x^B\) と \(x'^A = R^2/x^B\) で関係づけられる．
\end{proof}

\begin{remark}[記号 \(x'^A\) の意味について]
ここで \(x'^A\) は，観測者 \(B\) の主観座標系における \(A\) 方向の座標成分を意味する（論文 {[2]} 式 (6.5) 参照）．軸 \(A\) そのものではなく，\(B\) の座標系で \(A\) 方向を表す成分であることに注意．
\end{remark}

\subsection{観測量の役割の交換}

命題 4.1 から導かれる構造を整理する．背景空間の軸 \(A\) と軸 \(B\) について：

\begin{itemize}
\item \textbf{観測者 \(A\) から見ると}：軸 \(A\) は「投影中心軸」として直接アクセスできない外部の軸であり，軸 \(B\) は自分の主観空間内の座標 \(x^B\) である
\item \textbf{観測者 \(B\) から見ると}：軸 \(B\) は「投影中心軸」として直接アクセスできない外部の軸であり，軸 \(A\) は自分の主観空間内の座標 \(x'^A\) である
\end{itemize}

すなわち，背景空間における対等な二つの軸 \(A\) と \(B\) は，観測者の立場の選択によってその役割が交換される．論文 {[2]} 対称性 II（軸の対等性）が背景空間における軸の対等性を保証する一方，観測者の立場の選択によって軸の役割は非対称になる．

\begin{remark}[軸の対等性と役割の非対称性]
論文 {[2]} 対称性 II は背景空間における軸の対等性を述べる．本節の命題 4.1 は，主観空間の構成（中心投影）によってこの対等性が観測者の立場に依存する非対称性に転化することを示す．背景空間においては全ての軸が対等であるが，特定の軸を投影中心として選ぶと，その軸は観測者にとって「外部」となり，他の軸とは異なる役割を持つ．
\end{remark}

\begin{remark}[観測者の立場の任意性]
論文 {[2]} 対称性 II により，どの軸を投影中心として選んでも幾何学的構造は同一である．したがって，「どの観測者が正しい立場にあるか」という問いは幾何学的に意味を持たない．観測者 \(A\) と観測者 \(B\) の立場は完全に対等であり，両者の観測の差は観測者の選択そのものに起因する．
\end{remark}

\section{観測の相対性}

\subsection{主観空間の選択と観測量の関係}

前節の対比を一般化する．背景空間 \(\mathbb{R}^{n+1}\) から構成される複数の主観空間 \(\{H_A^+\}_{A=1}^{n+1}\) について，それぞれの主観空間における内部観測者は，同じ背景空間を共有しながら異なる観測量を取得する．

\begin{theorem}[観測の相対性]
背景空間 \(\mathbb{R}^{n+1}\) から構成される主観空間の族 \(\mathcal{H} = \{H_A^+\}_{A=1}^{n+1}\) について，次が成立する：

\textup{(i)} 各主観空間 \(H_A^+\) は構造的に同一の幾何学的性質を持つ（命題 2.1）

\textup{(ii)} 各主観空間 \(H_A^+\) における内部観測者 \(A\) は，計量 \(g_{\mu\nu}^{(A)}\) から構成される量を取得する（命題 3.1）

\textup{(iii)} 異なる主観空間 \(H_A^+\) と \(H_B^+\) における内部観測者 \(A\) と \(B\) は，同じ背景空間を共有しながら，それぞれの主観座標における異なる役割を持つ観測量を取得する（命題 4.1）

\textup{(iv)} 論文 {[2]} 対称性 II（軸の対等性）により，どの観測者の立場も幾何学的に対等であり，特定の観測者の立場を「正しい」とする幾何学的根拠は存在しない
\end{theorem}

\begin{proof}
(i)--(iii) は前節までの結果から直接従う．(iv) は論文 {[2]} 定理 4.1 による：軸の選択 \(A\) に対する写像・計量・曲率テンソルの構造は同一であり，どの軸を投影中心として選ぶかは幾何学的に区別できない．
\end{proof}

\subsection{観測の相対性の構造}

定理 5.1 が示すのは，中心投影の枠組みにおける観測量の相対性である．「観測量」は次の二段階で決まる：

\textbf{第一段階}：背景空間 \(\mathbb{R}^{n+1}\) における幾何学的構造．これは論文 {[2]} 対称性 II により全ての軸が対等で，普遍的に保存される．

\textbf{第二段階}：投影中心軸の選択（=観測者の立場の選択）．この選択により，背景空間の対等な軸のうち一つが「投影中心軸」として外部化され，残りの軸が観測者の主観空間内の座標として現れる．

異なる観測者は第二段階の選択が異なるため，それぞれ異なる観測量を取得する．しかし第一段階の背景空間は共通であり，論文 {[2]} 定理 6.1（主観座標系の変換可能性）により，異なる観測者の観測量は背景空間を経由した合成写像で相互変換可能である．

\begin{remark}[相対性の性格]
本論文で述べる「観測の相対性」は，観測者の主観空間の選択（=投影中心軸の選択）に依存する相対性である．これは内部観測者の慣性運動や加速運動の選択に依存する物理学的な相対性（ガリレイ相対性，特殊相対性，一般相対性など）とは異なる種類の相対性であり，より基底的な，主観空間そのものの選択についての相対性である．
\end{remark}

\section{結論}

本論文では，論文 {[2]} の対称性 II（軸の対等性）と対称性 IV（主観座標系の変換可能性）から自然に導かれる帰結として，次のことを示した：

\begin{enumerate}
\item 背景空間 \(\mathbb{R}^{n+1}\) から，投影中心軸の選択により \(n+1\) 通りの主観空間 \(\{H_A^+\}\) が同時に構成可能である（命題 2.1）

\item 各主観空間における内部観測者は，自分の主観空間の計量から構成される量のみを取得し，他の主観空間や背景空間の構造には直接アクセスできない（命題 3.1）

\item 異なる主観空間における内部観測者は，同じ背景空間を共有しながら，それぞれの主観座標における異なる役割を持つ観測量を取得する（命題 4.1）

\item 論文 {[2]} 対称性 II により，どの観測者の立場も幾何学的に対等であり，観測量の差は観測者の主観空間の選択に起因する（定理 5.1）
\end{enumerate}

これらの結果は，論文 {[2]} の対称性 II と対称性 IV の自然な帰結であり，新たな仮定を一切導入していない．本論文の主張はすべて純粋に幾何学的な命題であり，論文 {[1, 2]} の結果のみから論証される．

\section{本論文の主張範囲と留保}

本論文の主張範囲を明確にするため，以下の点を留保する．

\subsection{物理学的解釈について}

本論文は純粋に幾何学的な命題を扱う．中心投影，主観空間，対称性，これらはすべて微分幾何学の対象として定義され，物理学の特定の概念（時空，時間，重力場，電磁場，電荷，エネルギーなど）への対応は本論文では主張しない．読者が本論文の結果を物理学的に解釈することは可能であるが，そのような解釈は本論文の主張範囲を超え，別途の考察を要する．

特に，本論文で述べる「観測の相対性」は，幾何学的な構造についての主張であり，物理学における特定の理論（一般相対性理論，量子力学，宇宙論など）への対応を主張するものではない．

\subsection{主観空間の住民について}

本論文では「主観空間の内部観測者」「観測者 \(A\)」「観測者 \(B\)」といった表現を用いる．これらは中心投影によって構成される主観空間における内部観測者を指す比喩的な表現であり，物理的に実在する観測者の存在を主張するものではない．本論文は数学的構造として複数の主観空間を扱い，それぞれの主観空間における幾何学的性質を比較する．

「観測者」という表現は，論文 {[2]} \S2 用語表で定義された「内部観測者」（主観空間の内部で計量のみを用いて測定を行う仮想的な観測者）に従う．この定義は数学的なものであり，物理的な意識，知覚，主観性などの哲学的概念には踏み込まない．

\subsection{変換式の一般解について}

論文 {[2]} 定理 6.1（主観座標系の変換可能性）により，異なる二つの主観空間 \(H_A^+\) と \(H_B^+\) の主観座標の間には，球面 \(S^n(R)\) 上の点（神の目座標）を経由した合成写像
\[
T_{A \to B} = \Phi_B^{-1} \circ \Phi_A : \Pi_R^{(A)} \to \Pi_R^{(B)}
\]
が構造的に存在し，これは \(C^\infty\) 級の微分同相写像として正則である（共通領域 \(U_{AB} = H_A^+ \cap H_B^+\) に対応する主観座標の範囲において）．本論文で扱う「複数の主観空間における内部観測者が同じ背景空間を異なる形で観測する」という構造は，この合成写像の存在を前提としている．

ただし，本論文はこの合成写像の具体的な変換式を一般解として定式化していない．論文 {[2]} 命題 6.1 では2つの軸間の変換の明示的公式が与えられているが，任意の次元・任意の軸の組み合わせに対する閉形式の一般解の構成は，本論文の射程外であり，今後の研究課題である．

本論文の主張は，合成写像の存在（構造的証明）に基づく概念的な主張に留まる．変換式の具体的な構成は，個別の軸の組み合わせに対して論文 {[1, 2]} の中心投影の定式化から原理的には導出可能であるが，一般解としての閉形式の提示は別途の研究を要する．

\subsection{人間原理的表現について}

本論文では「観測者 \(A\)」「観測者 \(B\)」「観測者の立場」「観測者の選択」といった，人間原理的または神学的に解釈されうる表現を使用する．具体的には，中心投影の内部観測者を擬人的に記述し，異なる主観空間の内部観測者の取得情報を対比する．

これらの表現は，中心投影の幾何学的構造の理解を助けるための比喩的な表現であり，以下のいずれをも主張するものではない：

\begin{itemize}
\item 神の存在，あるいは神の視点の実在
\item 主観空間の内部観測者となる人間または他の存在の実在
\item 特定の主観空間における観測者の物理的存在
\item 観測という行為そのものの哲学的または意識的な性質
\item 人間原理（弱い人間原理，強い人間原理のいずれも）
\item 並行宇宙，マルチバース理論との対応
\item 量子力学の観測問題や意識の物理学
\end{itemize}

本論文の主張はすべて，中心投影という数学的構成における幾何学的対象の性質に関するものである．「観測者」「観測」「立場」といった表現は，読者の理解を助けるために導入された比喩であり，これらの表現から哲学的，神学的，または存在論的な結論を導くことは本論文の意図ではない．

本論文は物理学の基礎論文でも哲学の論文でもなく，微分幾何学における中心投影の対称性の構造を扱う論文である．

\section{結果の整理}

\begin{longtable}{p{4cm}p{7.5cm}p{3.5cm}}
\toprule
項目 & 内容 & 根拠 \\
\midrule
\endhead
\textbf{複数の主観空間の同時存在} & 背景空間 \(\mathbb{R}^{n+1}\) から \(n+1\) 通りの主観空間 \(\{H_A^+\}\) が構成可能 & 命題 2.1，論文 {[2]} 定理 4.1 \\
\addlinespace
\textbf{内部観測者の取得情報} & 観測者 \(A\) は主観空間 \(H_A^+\) の計量から構成される量を取得し，投影中心軸 \(A\) や他の主観空間にはアクセスしない & 命題 3.1，論文 {[2]} 定理 5.2 (i) \\
\addlinespace
\textbf{観測者の対比} & 観測者 \(A\) と観測者 \(B\) は同じ背景空間を共有しながら異なる観測量を取得する．背景空間の軸の役割は観測者の立場により交換される & 命題 4.1，論文 {[2]} 定理 6.1 \\
\addlinespace
\textbf{観測の相対性} & 観測量は観測者の主観空間の選択（投影中心軸の選択）に依存する．論文 {[2]} 対称性 II により観測者の立場は幾何学的に対等である & 定理 5.1，論文 {[2]} 対称性 II \\
\addlinespace
\textbf{変換式の存在（一般解は射程外）} & 異なる主観空間の間には背景空間を経由した \(C^\infty\) 級の正則変換が存在するが，本論文は閉形式の一般解を提示しない & \S7.3，論文 {[2]} 定理 6.1 \\
\bottomrule
\end{longtable}

\bigskip\hrule\bigskip

\section*{参考文献}

\noindent
{[1]} Kihara, N. (2026). 中心投影による4次元空間の幾何学的定式化. DOI: \href{https://doi.org/10.5281/zenodo.19427780}{10.5281/zenodo.19427780}.

\noindent
{[2]} Kihara, N. (2026). 中心投影の幾何学的対称性：多軸モデルの数学的基盤. DOI: \href{https://doi.org/10.5281/zenodo.19434932}{10.5281/zenodo.19434932}.

\noindent
{[3]} do Carmo, M. P., \textit{Riemannian Geometry}, Birkhaeuser, Boston, 1992, Chapter 4, 8.

\noindent
{[4]} Wald, R. M., \textit{General Relativity}, University of Chicago Press, Chicago, 1984, Appendix C, D.

\noindent
{[5]} Wolf, J. A., \textit{Spaces of Constant Curvature}, 6th ed., AMS Chelsea Publishing, Providence, 2011, Chapter 2.

\bigskip\hrule\bigskip

\begin{center}
\textit{本稿は純粋に幾何学的な命題の証明に限定される．物理的解釈，哲学的解釈，存在論的解釈は一切行わない．本論文の主張範囲は \S7 に明示的に限定される．}
\end{center}

\end{document}
